\section{Discusión de los resultados}\label{sec:discusion-resultados}

\subsection{Cumplimiento del problema planteado}

La solución resuelve el problema principal: conectar un catálogo técnico en OpenMetadata con una salida interoperable DCAT-AP-ES validable.\ct{Hecho: el tutor reorganizó este capítulo. 7.1-7.4 (cumplimiento, beneficios, limitaciones, amenazas) pasan a Resultados como ``Discusión de los resultados''; 7.5 (Big Data) y 7.6 (gobierno supervisado) se mueven a anexos; el trabajo futuro y las conclusiones se reestructuran en el capítulo ``Conclusiones y trabajos futuros''.} La plataforma no se limita a documentar un mapeo, sino que implementa un flujo ejecutable que descubre activos técnicos, permite curación funcional, aplica metadatos, exporta JSON-LD y valida el resultado. El Anexo~\ref{ap:demo-operativa} narra ese flujo de extremo a extremo desde la consola web, hasta dejar el archivo RDF listo para entregarlo a un harvester compatible.

La decisión de usar PostgreSQL de referencia en lugar de cosechar CKAN refuerza la defendibilidad del trabajo. CKAN habría permitido transformar metadatos ya publicados. PostgreSQL y OpenMetadata permiten demostrar gobierno desde activos técnicos, que es más coherente con las competencias de ingeniería, Big Data y cloud del máster.

\subsection{Beneficios obtenidos}

Los principales beneficios son:

\begin{itemize}
  \item reproducibilidad del entorno y del caso de uso;
  \item separación entre metadatos técnicos y curación funcional;
  \item núcleo Python único para CLI, scripts y web;
  \item validación formal mediante SHACL;
  \item trazabilidad entre SQL, OpenMetadata, hoja funcional, JSON-LD e informe de validación;
  \item facilidad de operación mediante consola web.
\end{itemize}

\subsection{Limitaciones}

La solución tiene limitaciones deliberadas:

\begin{itemize}
  \item una tabla SQL no equivale semánticamente a un \texttt{dcat:Dataset}; en el caso de uso se usa como aproximación técnica gobernable;
  \item se genera una única \texttt{Distribution} y un único \texttt{DataService} por dataset;
  \item la clasificación HVD es una hipótesis de validación dentro de la plataforma, no una calificación jurídica automática;
  \item la validación se centra en metadatos, no en profiling de contenido ni calidad fila a fila;
  \item no se publica automáticamente en datos.gob.es, CKAN ni data.europa.eu;
  \item el modelo OpenMetadata se mantiene mínimo y no cubre todos los campos opcionales de DCAT-AP-ES.
\end{itemize}

La primera limitación merece una precisión. Un \texttt{dcat:Dataset} es una abstracción editorial: una colección publicada y mantenida por un agente como unidad coherente, con independencia de su almacenamiento físico~\cite{w3cDcat3}. Una tabla \ac{SQL} es una estructura física, de modo que la correspondencia «una tabla es un dataset» es una decisión de modelado y no una equivalencia natural: un dataset podría corresponder igualmente a una vista, a una agregación o a la unión de varias tablas. En el caso de uso de validación se elige la tabla \texttt{gold} como unidad publicable por ser la representación más estable del dato refinado. En un escenario real, la hoja de gobierno permitiría asociar el dataset a la vista o agregación que mejor lo represente, sin alterar el núcleo.

Estas limitaciones no invalidan el trabajo: forman parte del alcance funcional declarado y delimitan la frontera entre lo que la plataforma demuestra de forma reproducible y lo que correspondería a una iniciativa de portal de datos abiertos a escala institucional, con presupuesto, equipo y compromisos de servicio diferenciados.

\subsection{Amenazas a la validez}

La principal amenaza a la validez externa es que el caso de uso trabaja con datos sintéticos y un número reducido de datasets publicables. La arquitectura, sin embargo, está diseñada para escalar a más tablas \texttt{gold} mediante la misma hoja funcional y el mismo workflow.

Conviene aclarar que la conformidad \ac{SHACL} «con warnings permitidos» no constituye una amenaza a la validez. En \ac{SHACL}, una violación con severidad \texttt{sh:Violation} indica el incumplimiento de una restricción obligatoria. Un \texttt{sh:Warning}, en cambio, señala únicamente que el perfil \emph{recomienda} cubrir propiedades adicionales, recomendadas u opcionales, que no son obligatorias~\cite{w3cShacl}. El catálogo exportado no presenta violaciones bloqueantes: los warnings solo invitan a enriquecer metadatos no exigidos por el perfil. Aumentar esa cobertura mejora la calidad editorial del catálogo, pero su ausencia no invalida la conformidad obligatoria demostrada.
